\documentclass{situationist-preprint}

\title{The Meaning of \enquote{Meaning}}
\msid{LNG-21-a}
\msfamily{LNG-21}
\msdiscipline{Philosophy of Language}
\mstype{Principal manuscript}
\mskeywords{meaning; ambiguity; significance; sense; meaning as use; polysemy}

\begin{document}
\maketitle

\begin{abstract}
We speak of the meaning of a word and of the meaning of a defeat, a gesture or a life. It is commonly assumed that \enquote{meaning} is ambiguous between linguistic sense and significance, and that philosophical theories of meaning concern only the first. We argue that the word is not ambiguous in this way, and that its meaning is fixed by the theory of LNG-01: the meaning of a thing is its role in a rule-governed practice, constrained by what that role contributes to how things are taken to be. Applied to words, this yields sense; applied to events and actions within practices and lives, it yields significance. The difference lies in the practices, not in the word. Standard tests for ambiguity confirm the verdict.
\end{abstract}

\section{The question}

\enquote{What does \enquote{vixen} mean?} asks for a sense. \enquote{What did the defeat at Sedan mean for the Second Empire?} asks for a significance. It is natural to conclude that \enquote{meaning} has two senses, and that the theory of meaning concerns only one. (Grice's distinction between natural and non-natural meaning, \citealp{Grice1957}, concerns the verb and is a different matter.) We argue that the word has one meaning, fixed by the account of LNG-01.

\section{The argument}

By LNG-01, the meaning of an expression is its use, the rules that govern it in a practice, constrained by its contribution to truth-conditions. Applied to \enquote{meaning} itself, the question is what rule governs the use of \enquote{the meaning of $x$}. The rule is this: the meaning of $x$ is what $x$'s role in a rule-governed practice makes so, or licenses taking to be so, within that practice. For a word, the practice is a language, and what the word's role contributes is its contribution to the truth-conditions of sentences; this is sense. For a defeat, a gesture or a vote, the practice is a political order, a ritual of greeting or an election, and what the item's role contributes is what it makes so within the practice: the end of a regime, an insult, a mandate. This is significance. One rule covers both, and the truth-conditional constraint applies in both cases, since \enquote{the defeat meant the end of the Empire} and \enquote{his silence meant consent} are true or false.

Two tests for ambiguity agree. First, conjunction reduction: \enquote{the word and the gesture had the same meaning}, said of \enquote{no} and a shake of the head, is not zeugmatic, as \enquote{she took her coat and her leave} is, with two senses of \enquote{take}. Second, the cross-linguistic test of \citet{Kripke1977}: if a word is ambiguous, other languages will tend to express its senses by different words. German \emph{Bedeutung} and French \emph{sens} each serve for both sense and significance, as \enquote{meaning} does.

\begin{objection}
\enquote{The meaning of life} has nothing to do with truth-conditions.
\end{objection}

\begin{reply}
The question asks what a life's place in the practices and projects of which it is part makes so, and its answers, such as \enquote{life has no meaning}, are true or false. Whether any such answer is true is a separate question; that the phrase applies the same rule as \enquote{the meaning of \enquote{vixen}} is our claim.
\end{reply}

\section{Conclusion}

\enquote{Meaning} is univocal. Its meaning is fixed by the theory of LNG-01, and the apparent distinction between sense and significance is a distinction between the practices in which items have roles.

\begin{thebibliography}{9}\small
\bibitem[Grice(1957)]{Grice1957} Grice, H.~P. (1957). Meaning. \emph{Philosophical Review}, 66(3), 377--388.
\bibitem[Kripke(1977)]{Kripke1977} Kripke, S. (1977). Speaker's reference and semantic reference. \emph{Midwest Studies in Philosophy}, 2, 255--276.
\end{thebibliography}

\end{document}
